\begin{abstract}

本文围绕NIPT检测中Y染色体浓度的时间动态规律，建立浓度建模、时点优化与异常判定的分析框架。

针对问题一，267名孕妇共1082条纵向观测数据的组内相关系数ICC=0.60，提示个体聚类效应不可忽略。本文建立\textbf{广义加性混合模型（GAMM）}，以B样条基函数拟合孕周的非线性效应，引入随机截距和随机斜率刻画个体差异。模型AIC为$-$5240.80，线性混合效应模型AIC=$-$5174.04，似然比检验\textbf{$\chi^2$=76.76}（$p<10^{-15}$），RMSE=\textbf{0.0101}，非线性成分的加入带来了实质性的拟合改善。Y染色体浓度随孕周非线性上升，BMI对浓度有负向稀释作用。

针对问题二，采用\textbf{数据驱动BMI分段与双风险函数优化}确定各组最佳检测时点。以BIC准则搜索最优断点，得到BMI=30和BMI=36两个分界，将孕妇分为三组。构造风险函数$R(t)$=0.4$R_{\text{delay}}(t)$+0.6$R_{\text{inaccuracy}}(t)$，同时权衡延迟风险与检测不准确风险。优化结果为G1组（BMI$\in$[20.7,30)）最优时点\textbf{11.7周}，G2组（BMI$\in$[30,36)）\textbf{11.8周}，G3组（BMI$\in$[36,46.9)）\textbf{13.5周}，高BMI组对误差敏感性明显更高。

针对问题三，建立\textbf{多因素GAMM并结合残差分解与Monte Carlo误差传播}分析扩展因素的影响。体重与BMI的共线性严重（VIF=351），残差分解后VIF降至\textbf{1.01}。检验发现年龄（$p$=0.309）和体重残差（$p$=0.685）均不显著，BMI单因素已足够。将孕妇细分为4组后最优时点分别为12.8、13.3、13.3、24.8周，500次MC模拟显示高BMI组95\%置信区间为[11.0,24.6]周，样本量不足是时点估计波动大的主要原因。

针对问题四，女性胎儿缺乏Y染色体标记，本文转而建立基于\textbf{多元马氏距离的异常检测方法}，选取Z13、Z18、Z21、ZX、BMI、X浓度共6维特征构建异常评分。该方法AUC=\textbf{0.574}，优于基线Z值阈值法（AUC=0.492），在第85百分位阈值处取得最佳F1=0.217，召回率26.9\%。检测能力受限于异常与正常样本在特征空间中的高度重叠。

本文的创新点在于，（1）GAMM捕捉孕周非线性效应，似然比检验量化改进幅度（$\chi^2$=76.76）；（2）双风险函数将NIPT时点选择转化为可优化的决策问题，定量权衡延迟风险与准确性风险；（3）残差分解消除体重与BMI共线性（VIF从351降至1.01），MC模拟量化高BMI组时点不确定性。

\end{abstract}

\keywords{NIPT检测；广义加性混合模型（GAMM）；BMI分段优化；Monte Carlo误差传播；马氏距离异常检测}
